\begin{musica}{Bella Ciao}{Popular italiana}{Tuna de Engenharia da Universidade do Porto}
Una mat\ac{Dm}{ti}na mi son svegliato
O bella ciao, bella ciao, bella \ac{D7}{ciao}, ciao, ciao,
Una mat\ac{Gm}{ti}na mi son svegli\ac{Dm}{ato}
E hotro \ac{A}{va}to\ac{A7}{}_ \ac{A}{l’i}n\ac{A7}{va}sor. \ac{Dm}{}__

O Partigiano portami via,
O bella ciao, bella ciao, bella ciao, ciao, ciao,
O Partigiano portami via
Che mi sento di morir.

E se io muoio da partigiano,
O bella ciao, bella ciao, bella ciao, ciao, ciao,
E se io muoio da partigiano
Tu me devi seppellir.

Seppellire sulla montagna,
O bella ciao, bella ciao, bella ciao, ciao, ciao,
Seppellire sulla montagna
Sotto l’ombra d’un bel fior.

E le genti che passeramo,
O bella ciao, bella ciao, bella ciao, ciao, ciao,
E le genti che passeramo
Mi diranno che biel fior.

Questo è il fiore del Partigiano,
O bella ciao, bella ciao, bella ciao, ciao, ciao,
Questo è il fiore del Partigiano
Morto per la Libertà.

Era rossa la sua bandiera,
O bella ciao, bella ciao, bella ciao, ciao, ciao,
Era rosso la sua bandiera
C’era scritto Libertà.

\posicoes{Dm: xx0231, D7: xx0212, Gm: 355333, A: x02220, A7: x02020}
\end{musica}
